% Part: counterfactuals
% Chapter: introduction
% Section: paradoxes-material

\documentclass[../../../include/open-logic-section]{subfiles}

\begin{document}

\olfileid{cnt}{int}{par}

\olsection{वास्तविक अभिव्यंजनाचे विरोधाभास}

इंग्रजीतील ‘‘जर \dots तर \dots’’
विधानांचे प्रतीकीकरण $!A \lif !B$
या रूपात करण्यावर टीका करणाऱ्यांपैकी
C.~I. Lewis हे सुरुवातीचे एक होते.
व्हाइटहेड आणि रसेल यांच्या
\emph{Principia Mathematica} मध्ये
वास्तविक अभिव्यंजनाचा वापर करताना
$\lif$ चा उच्चार ‘‘यावरून
तार्किकरीत्या निष्पन्न होते’’ असा
केला जात असे; त्यावर लुईस यांची
टीका होती. $\lif$ चा अर्थ
‘‘तार्किकरीत्या निष्पन्न होते’’
असा घेतला, तर कोणतेही असत्य
विधान~$p$ यावरून $p$ यावरून~$q$
निष्पन्न होते असे ठरेल, कारण
$p$ असत्य असताना
$p \lif (p \lif q)$ सत्य असते.
त्याचप्रमाणे कोणतेही सत्य
विधान~$q$ यावरून $p$ यावरून~$q$
निष्पन्न होते असे ठरेल, कारण
$q$ सत्य असताना
$q \lif (p \lif q)$ सत्य असते,
ही लुईस यांची रास्त तक्रार होती.

तर्कशास्त्रज्ञांना अर्थातच माहीत
आहे की \emph{तार्किक निष्पन्नता}
हा कारक नाही; तो !!{formula}s
किंवा विधानांमधील संबंध आहे.
म्हणून गोंधळ टाळण्यासाठी
$\lif$ ला ‘‘तार्किकरीत्या निष्पन्न
होते’’ असे वाचूच नये.\footnote{
  गणितज्ञ आणि संगणकशास्त्रज्ञ
  ‘‘$\lif$’’ चा उच्चार ‘‘यावरून
  निष्पन्न होते’’ असा अजूनही
  वारंवार करतात; पण लुईस यांनी
  दाखवलेला गोंधळ टाळण्यासाठी
  तत्त्वज्ञ ‘‘तरच’’ असा
  उच्चार करतात.} तसे
वाचत नसू, तर लुईस यांची
ती विशिष्ट चिंता उरत नाही:
$p$ हे एखादे असत्य इंग्रजी
विधान दर्शवत असले, तरी
$p$ यावरून $q$ तार्किकरीत्या
‘‘निष्पन्न’’ होत नाही.
$p \Entails q$ आहे का हे
ठरवण्यासाठी \emph{सर्व}
!!{valuation}s विचारात घ्याव्या
लागतात; योगायोगाने असत्य
ठरणारे विधान दर्शवण्यासाठी
$p$ वापरले, तरी
$p \Entails/ q$.

तरीही लुईस यांच्या पुढीलसारख्या
‘‘जर \dots तर\dots’’ विधानात
काहीतरी विचित्र वाटते:
\begin{quote}
चंद्र हिरव्या चीजचा बनलेला असेल,
तर $2+2=4$.
\end{quote}
आणि पुढील अनुमानांतही:
\begin{quote}
  चंद्र हिरव्या चीजचा बनलेला नाही.
  म्हणून चंद्र हिरव्या चीजचा
  बनलेला असेल, तर $2+2=4$.

  $2+2 = 4$. म्हणून चंद्र
  हिरव्या चीजचा बनलेला
  असेल, तर $2+2=4$.
\end{quote}
पण ‘‘जर \dots तर \dots’’
याचा अर्थ केवळ $\lif$
असा असेल, तर ते विधान
निर्विवाद सत्य आणि ही
अनुमाने निर्विवाद वैध ठरतील.

दुसरा प्रश्न $(!A \lif !B) \lor (!B \lif
!A)$ या सर्वतः सत्य सूत्राबद्दल
आहे. त्यावरून असे सुचेल की
वृत्तपत्रातून $S$ आणि $T$
ही कोणतीही दोन निश्चयार्थी
वाक्ये निवडली, तर ‘‘जर
$S$ तर $T$, किंवा जर
$T$ तर~$S$’’ हे विधान
सत्य असले पाहिजे.

\end{document}
